\section{Goal, boundaries and starting point}
I organize each sensor or IC as an independent project. Its folder contains all
application sources, portable driver sources, vendor MCU support, CubeMX settings,
Keil configuration and PlatformIO configuration. Copying a project does not require
the repository root or any sibling project. This first project uses DS18B20 on PB10
and STM32F107VCT6; the library itself does not depend on that processor.

The original copied application included LCD, relays, electrical measurement,
multiple temperature buses, Flash settings and build artifacts. Those functions
are removed from the new project. Only the development foundation and reviewed
portable protocol routines are retained. The raw template contains no sensor code.

\subsection{Hardware information versus verified measurements}
The copied configuration specifies an 8 MHz HSE crystal, PLL multiplier nine,
72 MHz core, 36 MHz APB1, 72 MHz APB2 and a 3.3 V design supply. The MCU has
256 KiB Flash and 64 KiB RAM\cite{stm32}. The user identifies PB10 as the connected
sensor pin. These facts define the firmware assumptions; I have not measured the
physical board, supply or waveform. Old PC7, relay and LCD connections are not used.

\subsection{What is delivered}
The project supports one addressed sensor or a caller-sized address array with
corresponding temperature and status arrays. Conversion scheduling is cooperative,
ROM discovery is incremental, and replacement preserves a slot when identity is
unambiguous. The portable implementation uses caller-owned static storage and no
heap. External sensor power is required. Parasite power, persistent mapping and
hardware-assisted transfer are future extensions, not implemented features.
